\section{Strutture algebriche}
\subsection{Gruppi, sottogruppi e semigruppi}
Indichiamo col termine \textbf{gruppo} un insieme \(G\) in cui viene definita
una operazione binaria \(\star\) che associa ogni coppia \(a,b \in G\) un risultato
\(c \in G\). Tale operazione deve soddisfare i seguenti assiomi:

\begin{itemize}

  \item 
  l'operazione gode della \textbf{proprietà associativa} \\
  \((a \star b) \star c = a \star (b \star c)\) \\

  \item 
  l'operazione è \textbf{interna} \\
  \( \forall a,b \in G \quad a \star b \in G\)

  \item
  esistenza dell'\textbf{elemento neutro} \\
  \( \forall a \in G \quad \exists!e \backepsilon a \star e = e \star a = a \) \\

  \item
  esistenza dell'\textbf{inverso} \\
  \( \forall a \in G \quad \exists!a' \backepsilon a \star a' = a' \star a = e \) \\

\end{itemize}

Inoltre, un gruppo viene definito \textbf{gruppo abeliano} se l'operazione \(\star\) 
gode anche della \textbf{proprietà commutativa}:

\[ \forall a,b \in G \quad a \star b = b \star a\]

Al contrario, se in un insieme \(H\) definiamo un'\textbf{operazione binaria interna} \(\odot\) 
che gode \underline{soltanto} della \textbf{propria associativa} allora indicheremo \(H\)
come un \textbf{semigruppo}.
\\
\\
Adesso che sappiamo cosa è un gruppo, parliamo dei \textbf{sottogruppi}. \\
Preso un gruppo \((G, \star)\), l'insieme \(H \subseteq G\) verrà indicato 
come \textbf{sottogruppo} di \(G\) se:

\begin{itemize}
  \item l'elemento neutro è presente sia nel gruppo "padre" che nel sottoinsieme \\ 
  \(e \in G, H\)

  \item chiusura dell'operazione \(\star\) \\
  \(\forall a,b \in H \quad a \star b \in H\)

  \item l'elemento inverso di un qualsiasi elemento di \(H\) sarà
  comunque interno all'insieme \(H\) \\
  \( \forall a \in H \quad a^{-1} \in H\)
  
\end{itemize}

\subsubsection{Gruppo quoziente}
Per poter parlare del \emph{grupo quoziente} è doveroso precisare cosa è
un \textbf{sottogruppo normale} (o \textbf{invariante}). \\
Preso un gruppo \(G\) e un suo sottogruppo \(H\) diremo che \(H\) è un 
sottogruppo normale se le sue \emph{classi laterali} destra e sinistra coincidono. \\
Cosa intendiamo per \textbf{classi laterali} allora? \\
Poichè \(G\) è un gruppo sappiamo che al suo interno è definita una 
operazione binaria \(\star\). Prendiamo quest'ultima e fissiamo un elemento 
\(g_k \in G\), la classe laterale \textbf{sinistra} del sottogruppo \(H\) saranno tutti
gli elementi del tipo:
\begin{align*}
  g_kH = \{ g_k \star h \} \quad \forall h \in H
\end{align*}
Analogamente, la classe laterale \textbf{destra} sarà l'insieme degli elementi
del tipo:
\begin{align*}
  Hg_k = \{ h \star g_k \} \quad \forall h \in H
\end{align*}
Continuando a verificare queste relazioni per ogni \(g \in G\)
otterremmo \textbf{tutte} le possibili classi laterali del sottogruppo \(H\),
e se ogni classe sinistra \(g_kH\) coincide con la rispettiva classe destra
\(Hg_k\) allora diremo che \(H\) è un \textbf{sottogruppo normale}.
\\
\\
Il \textbf{gruppo quoziente} sarà quindi l'insieme di tutte le classi 
laterali di un sottogruppo \(H\) e si indica come \(G/H\) ("\emph{G su H}").

\subsection{Anelli e sottoanelli}
Dopo aver definito \emph{gruppo} e \emph{sottogruppo} possiamo 
procedere con un'altra struttura algebrica: l'\textbf{anello}.
\\
\\
Definiamo in un insieme \(A\) le operazioni binarie interne di somma 
(\(+\)) e moltiplicazione (\(\cdot\)). \\
Se per ogni \(a,b,c \in A\):

\begin{itemize}
  
  \item \( (A, +) \) è un gruppo abeliano con elemento neutro \(0\)
  \item \( (A, \cdot) \) è un semigruppo
  \item La moltiplicazione è distributiva rispetto alla somma \\
  \( a \cdot (b + c) = ab + ac\) \\
  \( (a + b) \cdot c = ac + bc\)

\end{itemize}

Allora l'insieme \(A\) con le operazioni \( (+, \cdot) \) viene chiamato \textbf{anello}.\\
Inoltre, se la moltiplicazione gode anche della proprietà commutativa l'anello viene detto 
\textbf{anello commutativo}; se invece, la moltiplicazione ammette elemento neutro allora si dice
\textbf{anello unitario}. \\
Quando poi l'anello è di entrambi i tipi citati e non esistono divisori 
dello zero allora siamo di fronte
ad un \textbf{dominio d'integrità}. \\ 
Ma che vuol dire "non esistono divisori di zero"?
Vuol dire che se la moltiplicazione \(a \cdot b = 0\) allora sicuramente \(a = 0\) oppure \(b = 0\);
come già detto siamo "abituati" a pensare in un classico insieme \(\mathbb{R}\), quindi
una condizione di questo tipo per noi è "ovvia", ma non è sempre così.
\\
\\
Prendiamo come esempio l'insieme numerico \(\mathbb{Z}_4\). \\
Il numero al pedice ci indica che siamo nel "classico" insieme \(\mathbb{Z}\) 
ma al suo interno troviamo soltanto i possibili resti di una qualsiasi divisione per \(4\).
Una divisione \(n : a\) (dove \(n\) è il nostro pedice) avrà sempre come risultato un numero con resto compreso tra 
\(0 \le r \le n - 1\), o più semplicemente, tutti gli elementi di \(\mathbb{Z}_4\)
saranno "in modulo \(4\)". \\
L'insieme \(\mathbb{Z}_4\) sarà quindi composto così:
\[ \mathbb{Z}_4 = \{ 0, 1, 2, 3\} \]
Più avanti vedremo come la notazione più corretta sarebbe:
\[ \mathbb{Z}_4 = \{ [0], [1], [2], [3]\} \]
Ma non è importante per il momento. \\
Come detto prima non è sempre "ovvio" che: 
\[ a \cdot b = 0 \iff a = 0 \lor b = 0\]

Poichè siamo in \(\mathbb{Z}_4\) una moltiplicazione come
\(2 \cdot 2\) non avrà lo stesso risultato dell'insieme \(\mathbb{R}\); 
sarà infatti:
\[ 2 \cdot 2 \equiv 0 \pmod{4}\]
E in questo caso \(a \cdot b = 0\) con \(a, b \neq 0\)
\\
\\
Dopo aver approfondito gli anelli possiamo ora parlare dei \textbf{sottoanelli}.
Preso un anello \(A\), dove quindi sono definite le operazioni binarie interne 
di somma e prodotto, e che queste godono delle proprietà già citate, diremo che 
\(S\) è un \textbf{sottoanello} di \(A\) se:

\begin{itemize}
  \item l'elemento neutro \(0\) appartiene sia all'anello "grande" \(A\)
  che ad \(S\)

  \item chiusura rispetto alla sottrazione \\
  \( \forall a,b \in S \quad a - b \in S \)

  \item chiusura rispetto alla moltiplicazione \\
  \( \forall a,b \in S \quad a \cdot b \in S \)
\end{itemize}

\subsection{Ideali}
Questa tipologia di struttura algebrica è simile a quella di un sottoanello poichè
anche gli \textbf{ideali} sono sottoinsiemi di anelli, ma seguono le proprie regole. \\
Prendiamo la tupla \( (A, +, \cdot) \), dove \(A\) è un anello. Preso un sottoinsieme 
\(I \subseteq A\) diremo che è un \textbf{ideale destro} se:

\begin{itemize}
  \item \( (I, +) \) è un sottogruppo di \( (A, +) \)
  \item \( \forall i \in I \quad i \cdot a \in I \)
\end{itemize}

al contrario, diremo che \(I\) è un \textbf{ideale sinistro} se:

\begin{itemize}
  \item \( (I, +) \) è un sottogruppo di \( (A, +) \)
  \item \( \forall i \in I \quad a \cdot i \in I \)
\end{itemize}

Nel caso in cui \(A\) è un \textbf{anello commutativo} non faremo nessuna distinzione
tra ideale sinistro o destro e diremo solo che \(I\) è un \textbf{ideale}.
\\
\\
Useremo invece il termine \textbf{ideale proprio} quando un ideale \(I\) è un
sottoinsieme proprio di \(A\) (\( I \subset A\)); ne segue che un anello \(A\) può avere
infiniti ideali propri \(I_n\), ma tra tutti questi solo uno è l'\textbf{ideale massimale},
overo un ideale proprio che non è contenuto in nessun altro ideale proprio. \\
Questo tipo di ideale si chiama così perché è il "soffitto" di tutti gli ideali \(I_n\):
oltre lui c'è soltanto l'anello padre \(A\).
\\
\\
Infine chiameremo \textbf{ideale primo} quel tipo di ideale proprio in cui vale:
\[ \forall a,b \in A \quad a \cdot b \in I \implies a \in I \lor b \in I\]

\subsection{Campi}
Infine abbiamo la struttura più semplice ma che voglio comunque appuntare
per correttezza: il \textbf{campo}. \\
Presa una tupla \( (K, +, \cdot) \) dire che \(K\) è un campo se:
\begin{itemize}
  \item \( (K, +) \) è un gruppo abeliano con elemento neutro \(0\)
  \item \( (K \setminus \{0\}, \cdot) \) è un gruppo abeliano con elemento neutro \(0\)
  \item l'operazione di moltiplicazione gode della proprietà distributiva rispetto all'addizione
\end{itemize}
Quello di cui a noi davvero importa però sono i \textbf{campi finiti}, ovvero
campi il cui numero di elementi (o \emph{cardinalità}) è un numero finito. 
Questo poichè ogni campo finito ha \(p^n\) elementi, dove \(p\) è un numero primo
e \(n\) è un numero naturale \( \ge 1\). \\
La cosa importante da notare è che se noi costruissimo due campi finiti "diversi"
(ovvero con elementi diversi) finiremo comunque per costruire due campi che sono 
\emph{strutturalmente uguali}, potremo metterci qualsiasi nome vogliamo,
potremo costruirli in maniera diversa, in ogni caso finiremo con due campi 
identici.